\documentclass[11pt,a4paper]{article}

\usepackage[utf8]{inputenc}
\usepackage[T1]{fontenc}
\usepackage{lmodern}   % scalable fonts (needed for microtype font expansion)
\usepackage{amsmath,amssymb,amsfonts}
\usepackage{mathtools}
\usepackage{bm}
\usepackage{geometry}
\usepackage{booktabs}
\usepackage{longtable}
\usepackage{array}
\usepackage{xcolor}
\usepackage{listings}
\usepackage{hyperref}
\usepackage{parskip}
\usepackage{enumitem}
\usepackage{microtype}

\geometry{margin=2.5cm}

\hypersetup{
  colorlinks=true,
  linkcolor=blue!60!black,
  urlcolor=blue!60!black,
  citecolor=blue!60!black
}

\lstset{
  basicstyle=\ttfamily\small,
  backgroundcolor=\color{gray!10},
  frame=single,
  framerule=0.4pt,
  breaklines=true,
  columns=fullflexible,
}

% Short-hands
\newcommand{\vp}{v_\perp}
\newcommand{\vpa}{v_\parallel}
\newcommand{\vpmin}{v_{\perp,\min}}
\newcommand{\vpmax}{v_{\perp,\max}}
\newcommand{\vpamin}{v_{\parallel,\min}}
\newcommand{\vpamax}{v_{\parallel,\max}}
\newcommand{\fvdf}{f(\vp,\vpa)}
\newcommand{\dt}{\Delta t}
\newcommand{\Lcal}{\mathcal{L}}
\newcommand{\namvar}[1]{\texttt{\textbf{#1}}}

% ============================================================
\title{\textbf{FP2D\_QLRF\_NL}\\[0.4em]
       \large User Manual\\[0.4em]
       \normalsize 2D Fokker--Planck solver for ion velocity distribution
       functions in cylindrical velocity space}
\author{F.~Louche\\LPP-ERM/KMS}
\date{Version 2.3 -- July 2026}
% ============================================================

\begin{document}

\maketitle
\tableofcontents
\newpage

% ============================================================
\section{Introduction}
% ============================================================

\texttt{FP2D\_QLRF\_NL} solves the two-dimensional Fokker--Planck (FP) equation
for the velocity distribution function (VDF) $f(\vp,\vpa,t)$ of a minority ion
species in a magnetised plasma.  The two velocity coordinates are the speed
perpendicular to the magnetic field $\vp \ge 0$ and the speed parallel to the
field $\vpa \in (-\infty,+\infty)$.  Azimuthal symmetry about the magnetic
field direction is assumed throughout.

The code accounts for three physical operators:
\begin{enumerate}
  \item \textbf{Coulomb collisions} with one or more Maxwellian background
        species (electrons and bulk ions).
  \item \textbf{Quasi-linear radio-frequency (RF) heating} via the
        quasi-linear diffusion operator, evaluated using complex Bessel
        functions.
  \item \textbf{Self-collisions} (collisions of minority ions with
        themselves), with five levels of approximation selectable at runtime.
\end{enumerate}
An optional NBI \textbf{beam source} term and an associated particle
\textbf{loss term} are also included.

The code can compute either a \emph{steady-state} solution (direct solve) or
a \emph{time-dependent} evolution (implicit time integration).  Spatial
discretization uses a 7-point Fornberg finite-difference stencil.  The
resulting sparse linear systems are solved with Intel MKL PARDISO.


% ============================================================
\section{Physical Model}
% ============================================================

\subsection{The Fokker--Planck Equation}

The VDF $f = \fvdf$ evolves according to
\begin{equation}
\label{eq:FP}
  \frac{\partial f}{\partial t} =
    A\,f
  + B\,\frac{\partial f}{\partial \vp}
  + C\,\frac{\partial f}{\partial \vpa}
  + D\,\frac{\partial^2 f}{\partial \vp^2}
  + E\,\frac{\partial^2 f}{\partial \vp\,\partial \vpa}
  + F\,\frac{\partial^2 f}{\partial \vpa^2}
  + S(\vp,\vpa),
\end{equation}
where the six coefficient arrays $A$--$F$ (called \texttt{all00}, \texttt{all10},
\texttt{all01}, \texttt{all20}, \texttt{all11}, \texttt{all02} in the code) are
assembled as the sum of three independently computed contributions:
\begin{equation}
  X = X_{\mathrm{coll}} + X_{\mathrm{RF}} + X_{\mathrm{SC}},
  \qquad X \in \{A,B,C,D,E,F\},
\end{equation}
for Coulomb collisions, RF heating, and self-collisions, respectively.
The source term $S$ collects the NBI beam injection and the particle loss.

In steady-state mode ($\partial f/\partial t = 0$) the right-hand side of
\eqref{eq:FP} is set to zero and the resulting linear system is solved
directly.

\subsection{Coulomb Collision Operator}

The Coulomb collision operator is the linearized Fokker--Planck operator
describing the interaction of the minority ion species~$a$ (charge $Z_a$,
mass $m_a$, thermal velocity $v_a$) with a Maxwellian background species~$b$
(charge $Z_b$, mass $m_b$, thermal velocity $v_{Tb}$, density $n_b$).
The coefficient arrays $(\texttt{colin00}, \ldots, \texttt{colin02})$ for
each background species are computed by the subroutine \texttt{cblin} and
summed over all \texttt{nbulk} background species.

The collision frequency scale factor is
\begin{equation}
  \Gamma_b = \frac{n_b\,Z_a^2\,Z_b^2\,e^4\,\ln\Lambda_{ab}}
                  {4\pi\,\varepsilon_0^2\,m_a^2},
\end{equation}
where $\ln\Lambda_{ab}$ is the Coulomb logarithm evaluated at each time step
(see Section~\ref{sec:coulomb_log}).

\subsubsection{Coulomb Logarithm}\label{sec:coulomb_log}

Two distinct formulas are used:

\paragraph{Ion--electron:}
\begin{equation}
  \ln\Lambda_{ae} = 24 - \ln\!\left(\frac{n_e^{1/2}}{T_e}\right),
\end{equation}
with $T_e$ in eV and $n_e$ in m$^{-3}$.

\paragraph{Ion--ion:}
\begin{equation}
  \ln\Lambda_{ab} = 23 - \ln\!\left[
    Z_a^2 Z_b^2\,\frac{m_a+m_b}{m_a\,T_b + m_b\,T_a}\,
    \left(\frac{n_a\,Z_a^2}{T_a}+\frac{n_b\,Z_b^2}{T_b}\right)^{1/2}
  \right].
\end{equation}
In time-dependent runs the Coulomb logarithm is recomputed at each time step
using the current effective temperature of the minority ion species.

\subsection{Quasi-Linear RF Operator}

When radio-frequency heating is active (\texttt{irf} $= -1$), the
quasi-linear diffusion operator is added to the coefficient arrays.  The RF
operator describes the resonant power exchange between the minority ions and
an electromagnetic wave characterised by:
\begin{itemize}
  \item Wave frequency $\omega = 2\pi f$ (\texttt{frek} in Hz).
  \item Parallel wave vector $k_\parallel$ (\texttt{kpar} in rad/m).
  \item Perpendicular wave vector $k_\perp$ (complex, \texttt{kperp}
        in rad/m).
  \item Left- and right-hand polarisation components $E_+$, $E_-$
        (\texttt{eplus}, \texttt{emin}).
  \item Static magnetic field $B_0$ (\texttt{b0} in T).
  \item Harmonic number $n$ (\texttt{nharm}).
\end{itemize}

The quasi-linear diffusion coefficients are computed in the subroutine
\texttt{qlrfterm} using complex Bessel functions $J_n(k_\perp \vp / \omega_c)$
and their derivatives, evaluated at each grid point.  The resonance condition
is $\omega - n\,\omega_c - k_\parallel\,\vpa = 0$, broadened by a Gaussian
of width \texttt{delta\_RF} in $\vpa$.

\subsection{Self-Collision Operator}

Four options are available for the treatment of self-collisions, selected by
the flag \texttt{isc}:

\paragraph{\texttt{isc = 0} -- No self-collisions.}
The self-collision coefficients are zero.  The FP equation is linear.

\paragraph{\texttt{isc = 1} -- Maxwellian background at fixed temperature.}
The self-collision operator is computed using the Rosenbluth potentials for a
Maxwellian distribution at the Stix effective temperature $T_\mathrm{eff}$
(computed once at the start of the simulation and kept constant).  This adds
a fixed linear operator to the coefficient arrays.

\paragraph{\texttt{isc = 2} -- Maxwellian background at varying effective temperature.}
Same as \texttt{isc = 1}, but the background temperature and Coulomb
logarithm are updated at each time step to track the evolving effective
temperature $T_\mathrm{eff}(t) = (2T_\perp + T_\parallel)/3$.  Requires a
time-dependent run.

\paragraph{\texttt{isc = 3} -- Maxwellian background at the density-characteristic
temperature.}
Same as \texttt{isc = 2}, but the background temperature is the
\emph{density-characteristic} (cold-bulk) temperature $T_n(t)$ instead of the
energy-weighted $T_\mathrm{eff}(t)$.  $T_n$ is obtained from a
phase-space-weighted least-squares fit of $\ln f$ versus $v^2$ restricted to the
dense thermal core --- only cells with $f > \texttt{core\_frac}\cdot\max f$ enter
the fit --- which returns the temperature exactly for a Maxwellian.  Being tied to
the cold bulk rather than to the supra-thermal tail, this background better
represents the population that dominates the collisional drag.  The core fraction
\texttt{core\_frac} (default $3.8\times10^{-3}$) is tunable; the appropriate value
depends on the shape of the distribution, so it may need re-calibrating against the
fully nonlinear operator for a new scenario.  Requires a time-dependent run.

\paragraph{\texttt{isc = -1} -- Fully nonlinear self-collisions.}
The complete nonlinear Landau/Fokker--Planck self-collision operator
\begin{equation}
  \left(\frac{\partial f}{\partial t}\right)_{\mathrm{SC}}
  = \Gamma_{aa}\,\frac{\partial}{\partial v_i}
    \left[ H_{ij}[f]\,\frac{\partial f}{\partial v_j}
         - G_i[f]\,f \right],
\end{equation}
is used, where $H_{ij}$ and $G_i$ are the Rosenbluth potentials built from
the \emph{current} distribution $f(t)$.  The Rosenbluth potentials are
evaluated via a precomputed phi-distance kernel \texttt{sum\_phi} using
Gauss--Legendre quadrature (subroutine \texttt{distance\_v\_gauss\_legendre}).
This kernel depends only on the velocity grid and is computed once and saved
to disk (\texttt{sum\_phi.dat}); it can be reloaded for subsequent runs on the
same grid.  The operator is rebuilt at every time step, making the problem
nonlinear.  Requires a time-dependent run.

\subsection{Beam Source and Particle Losses}\label{sec:beam}

When \texttt{isource} $= -1$, a neutral-beam injection (NBI) source is added:
\begin{equation}
  S(\vp,\vpa) = \frac{n_0}{\tau_s}\,
    \frac{\delta_G(\vp - \vp^\mathrm{beam})\,\delta_G(\vpa - \vpa^\mathrm{beam})}
         {\mathcal{N}},
\end{equation}
where $\delta_G(x) = e^{-(x/\sigma)^2}/(\sqrt{\pi}\,\sigma)$ is a Gaussian
approximation to the Dirac delta, $\vp^\mathrm{beam} = v_\mathrm{beam}\sin\alpha$,
$\vpa^\mathrm{beam} = v_\mathrm{beam}\cos\alpha$ with $\alpha$ the injection
angle, and $\mathcal{N}$ is the normalization factor
\begin{equation}
  \mathcal{N} = \sqrt{\pi}\,\sigma_\perp\,e^{-(\vp^\mathrm{beam}/\sigma_\perp)^2}
               + \pi\,\vp^\mathrm{beam}\,\mathrm{erfc}(-\vp^\mathrm{beam}/\sigma_\perp),
\end{equation}
which reduces to $2\pi\,\vp^\mathrm{beam}$ for $\vp^\mathrm{beam} \gg \sigma_\perp$
and to $\sqrt{\pi}\,\sigma_\perp$ for a purely parallel beam ($\vp^\mathrm{beam}=0$).

The loss term is $-f/\tau_s$ (isotropic particle losses with slowing-down time $\tau_s$).


% ============================================================
\section{Numerical Methods}
% ============================================================

\subsection{Velocity-Space Grid}

The computational domain is $[\vpmin,\,\vpmax] \times [\vpamin,\,\vpamax]$,
discretized on a $N_\perp \times N_\parallel$ grid with 1-based indices
$(i,j)$ where $i \in [1,N_\perp]$ and $j \in [1,N_\parallel]$.

Three grid spacing modes are available, controlled by \texttt{ising}:

\begin{description}
  \item[\texttt{ising = 0}] Uniform spacing in both directions.
        $\Delta\vp = (\vpmax-\vpmin)/(N_\perp-1)$,
        $\Delta\vpa = (\vpamax-\vpamin)/(N_\parallel-1)$.

  \item[\texttt{ising = -1}] Two-domain grid in $\vp$.  A finer mesh of
        \texttt{nsing} points is used for $\vp \le \texttt{vbound}$, and a
        coarser uniform mesh for $\vp > \texttt{vbound}$.  This provides
        higher resolution near the magnetic axis where gradients are typically
        steeper.

  \item[\texttt{ising = +1}] Quadratic spacing in $\vp$, providing higher
        resolution near the axis $\vp = 0$.  The \texttt{fd\_stencil\_2d}
        routine uses the Fornberg algorithm (\texttt{fornberg\_weights}) to
        compute exact finite-difference weights for the local non-uniform
        spacing, so this mode is fully supported.
\end{description}

The cylindrical Jacobian is $J = 2\pi\,\vp$, so that integrals over velocity
space are $\int\!\int f\,J\,\mathrm{d}\vp\,\mathrm{d}\vpa$.

The global 1D degree-of-freedom (DOF) index is
\begin{equation}
  \mathrm{ix}(i,j) = (i-1)\,N_\parallel + j,
  \quad i \in [1,N_\perp],\; j \in [1,N_\parallel],
\end{equation}
giving $N_\mathrm{dof} = N_\perp \times N_\parallel$ total DOFs.

\subsection{Finite-Difference Discretization}

Spatial derivatives are discretized using the \textbf{7-point Fornberg
stencil} (subroutine \texttt{fd\_stencil\_2d}), which is a fourth-order
accurate scheme.  The stencil uses 4 points in each direction: at the current
node and its 3 neighbours.  For a non-uniform $\vp$ grid, the Fornberg
algorithm computes the exact derivative weights for the local grid spacing via
subroutine \texttt{fornberg\_weights}.

The stencil returns at most 49 entries per row (7~points in $\vp$ by 7 in
$\vpa$), although boundary rows have fewer entries.

\subsubsection{Boundary Conditions}

\begin{description}
  \item[Axis ($i=1$, $\vp=\vpmin$):] Neumann condition $\partial f/\partial\vp = 0$,
        enforcing axial symmetry.

  \item[Outer boundary ($i=N_\perp$):] Dirichlet condition $f = 0$.

  \item[Parallel boundaries ($j=1$ and $j=N_\parallel$):] Dirichlet condition $f = 0$.
\end{description}

\subsection{Advection Stabilization at High $\vp$ (\texttt{i\_upwind})}
\label{sec:upwind}

The perpendicular drag coefficient $B$ (collisional friction) grows with $\vp$
while the perpendicular diffusion $D$ falls off on the supra-thermal tail, so the
cell-P\'eclet number
\begin{equation}
  \mathrm{Pe}_\perp = \frac{|B|\,\Delta\vp}{D}
\end{equation}
increases steeply toward the outer boundary and can reach $\mathrm{Pe}_\perp \sim
40$ at $\vp = \vpmax$ for high-energy (e.g.\ second-harmonic ICRF) scenarios.
Beyond $\mathrm{Pe}_\perp \gtrsim 2$ a central discretization of the advection
term is no longer dissipative; at the outer Dirichlet wall, where the Fornberg
stencil is one-sided, the effective perpendicular diffusion of the discrete
operator can even become \emph{negative}.  A modified-equation analysis of the
Dirichlet-truncated wall row gives, in the uniform-spacing limit,
\begin{equation}
  D_{\mathrm{eff}} \;=\; \frac{223}{360}\,D \;-\; \frac{1}{12}\,|B|\,\Delta\vp ,
\end{equation}
i.e.\ the discrete diffusion turns negative exactly when the wall-cell P\'eclet
number exceeds $223/30 \approx 7.4$ (the deleted wall node biases the drag stencil
downwind; high order shrinks the classic two-point downwind coefficient from
$\tfrac12$ to $\tfrac{1}{12}$ but cannot change its sign).  The assembled operator $\Lcal$ then
acquires a small set of eigenvalues with positive real part, localized in a thin
layer against the $\vp = \vpmax$ boundary.  Under Crank--Nicolson
(\texttt{icn = -1}) these modes are amplified and the run eventually blows up,
whereas the fully-implicit scheme damps them.  The instability is governed by
$\vpmax$ (the wall velocity) and is essentially independent of the RF term and of
the grid resolution.

Setting \texttt{i\_upwind = -1} cures this at the discretization level: in every
cell with $\mathrm{Pe}_\perp > 2$ the drag term $B\,\partial f/\partial\vp$ is
evaluated with a first-order \emph{upwind} one-sided difference (forward for
$B>0$, backward for $B<0$, i.e.\ the side that makes the semi-discrete eigenvalue
dissipative), while the diffusion term keeps its central Fornberg weights.  In the
modified equation the upwind drag moment $+|B|\,\Delta\vp/2$ replaces
$-|B|\,\Delta\vp/12$, so $D_{\mathrm{eff}} > 0$ at \emph{any} P\'eclet number; the
$\mathrm{Pe}_\perp > 2$ switching threshold is therefore conservative with respect
to the $\approx 7.4$ instability threshold above.  The
sparsity pattern is unchanged (the diffusion stencil already fills every
same-column entry), so the phased-PARDISO machinery is untouched.  This removes
all positive-real-part modes with a bulk-physics change of at most a few percent,
because the upwinding acts only on the near-empty tail.  The default
\texttt{i\_upwind = 0} keeps the pure central scheme; running fully implicit
(\texttt{icn = 0}) is an alternative stabilization.  Note that first-order
upwinding adds some numerical diffusion on the tail; a flux-limited or
exponentially-fitted variant would reduce this if higher tail accuracy is needed.

\subsection{Sparse Matrix Storage}

All matrices are stored in \textbf{1-based Compressed Sparse Row (CSR)} format.
For a matrix of size $N_\mathrm{dof}$:
\begin{itemize}
  \item \texttt{ia(1:$N_\mathrm{dof}$+1)}: row pointer array.
  \item \texttt{ja(1:nnz)}: column indices.
  \item \texttt{aa(1:nnz)}: non-zero values.
\end{itemize}
Entries within each row are sorted by column index before passing to PARDISO.
The maximum number of non-zeros per row is 49 (7$\times$7 stencil), so
$\texttt{nnz\_max} = 49\,N_\mathrm{dof}$.

\subsection{Direct Solver: Intel MKL PARDISO}

All linear systems are solved with \textbf{Intel MKL PARDISO} using matrix
type \texttt{mtype = 11} (real non-symmetric general).  The module
\texttt{pardiso\_solver} exposes four phased entry points that exploit
PARDISO's factorization phases:

\begin{center}
\begin{tabular}{ll}
\toprule
Entry point & PARDISO phases \\
\midrule
\texttt{pardiso\_solve\_steady}  & 11 + 22 + 33 (full solve, then release) \\
\texttt{pardiso\_solve\_init}    & 11 + 22 (symbolic + numeric factorization) \\
\texttt{pardiso\_solve\_step}    & 33 only (back-substitution) \\
\texttt{pardiso\_solve\_finalize}& release internal memory \\
\bottomrule
\end{tabular}
\end{center}

This phased approach is critical for performance in the time-dependent solver:
when the LHS matrix is constant (linear problem), the factorization is done
once and only the back-substitution is repeated at each time step.


% ============================================================
\section{Algorithms}
% ============================================================

\subsection{Steady-State Solver}

When \texttt{ntimes(1) = 0}, the code solves the steady-state FP equation
\begin{equation}
  \Lcal[f] = -S,
\end{equation}
where $\Lcal$ is the FP differential operator assembled from the coefficient
arrays \texttt{all00}, \ldots, \texttt{all02}.  The two-pass assembly (first
count non-zeros, then fill values) builds the CSR matrix, which is then
passed to \texttt{pardiso\_solve\_steady} for a direct solve (subroutine
\texttt{FP\_steady\_state} in module \texttt{mod\_linear}).

\subsection{Time-Dependent Solver -- Linear Problem (\texttt{timefp\_7pt})}

Used when \texttt{ntimes(1) > 0} and $\texttt{isc} \ne -1$.

\subsubsection{Time Discretization}

A $\theta$-method is used:
\begin{equation}
  \bigl(I - \theta\,\dt\,\Lcal\bigr)\,f^{n+1}
  = \bigl(I + (1-\theta)\,\dt\,\Lcal\bigr)\,f^n + \dt\,S.
\end{equation}

\begin{center}
\begin{tabular}{lll}
\toprule
\texttt{icn} & $\theta$ & Scheme \\
\midrule
$-1$  & $0.5$  & Crank--Nicolson \\
$+1$  & $0.75$ & Intermediate \\
else  & $1.0$  & Fully implicit \\
\bottomrule
\end{tabular}
\end{center}

\subsubsection{Algorithm}

\begin{enumerate}
  \item \textbf{Assemble} the FP operator $\Lcal$ in CSR format (two-pass loop).
  \item \textbf{Build LHS}: $M_\mathrm{lhs} = I - \theta\,\dt\,\Lcal$.  The sparsity
        pattern is identical to $\Lcal$ since the diagonal already exists.
  \item \textbf{Factorize} $M_\mathrm{lhs}$ once (PARDISO phases 11+22).
  \item \textbf{Time loop} for $n = 1, \ldots, N_\mathrm{times}$:
        \begin{enumerate}
          \item Recompute the Coulomb logarithm and, for
                \texttt{isc} $\in \{1,2,3\}$, the Maxwellian self-collision
                background (at the fixed Stix, $T_\mathrm{eff}(t)$, or $T_n(t)$
                temperature, respectively) from the current distribution, then
                rebuild the FP operator (allowing a time-varying Coulomb logarithm
                and self-collision term).
          \item Rebuild $M_\mathrm{lhs}$ with the updated operator.
          \item Compute the RHS:
                $b = \bigl(I + (1-\theta)\,\dt\,\Lcal\bigr)\,f^n + \dt\,S$
                via a hand-coded sparse matrix--vector product.
          \item Enforce zero RHS on boundary DOFs.
          \item Solve $M_\mathrm{lhs}\,f^{n+1} = b$ (PARDISO phase 33 only).
          \item Compute diagnostics (density, energy, power, momentum).
        \end{enumerate}
  \item \textbf{Finalize} PARDISO.
  \item If \texttt{isource = 0}: rescale $f \leftarrow f \cdot n_0/\int f\,\mathrm{d}^3v$
        to conserve total particle number.
\end{enumerate}

\subsection{Time-Dependent Solver -- Nonlinear Self-Collisions (\texttt{timefp\_7pt\_nl})}

Used when \texttt{ntimes(1) > 0} and \texttt{isc = -1}.  The LHS matrix now
changes at every time step because the self-collision operator depends on
$f^n$.

\subsubsection{Phi-Distance Kernel Precomputation}

Before entering the time loop, the phi-distance kernel
$\Phi_{ij} = \phi(|\bm{v}_i - \bm{v}_j|)$ (Gauss--Legendre quadrature,
$O(N_\mathrm{dof}^2)$) is cached on disk in a per-case file
\texttt{sum\_phi-CASENAME.dat} (or \texttt{sum\_phi.dat} when \texttt{casename}
is empty). Since the kernel is grid-dependent, the file header stores a grid
signature (\texttt{nperp}, \texttt{npar}, \texttt{vperp\_min/max},
\texttt{vpar\_min/max}). Caching is fully self-managing --- no user flag is
required:

\begin{center}
\begin{tabular}{ll}
\toprule
Cache file state & Action \\
\midrule
Missing                                  & Compute kernel and save \\
Present, signature matches               & Load kernel (skip the quadrature) \\
Present, signature differs or truncated  & Recompute and overwrite \\
\bottomrule
\end{tabular}
\end{center}

Because the signature is verified automatically, a stale kernel from a
different velocity grid can never be silently reused.

\subsubsection{Time-Stepping Scheme}
\label{sec:nl_theta}

The $\theta$ parameter has an additional option for the NL solver:

\begin{center}
\begin{tabular}{lll}
\toprule
\texttt{icn} & $\theta$ & Comment \\
\midrule
$-1$ & $0.5$ & Crank--Nicolson (may be unstable at high $\vp$; see Sec.~\ref{sec:upwind}) \\
$+1$ & $0.75$ & Intermediate \\
else & $1.0$ & Fully implicit (recommended) \\
\bottomrule
\end{tabular}
\end{center}

\subsubsection{Algorithm}

\begin{enumerate}
  \item \textbf{Build sparsity pattern} of $\Lcal$ using the linear
        coefficients only (pattern is fixed; NL SC uses the same stencil).
  \item \textbf{Factorize symbolically} (PARDISO phase 11 only).
  \item \textbf{Time loop} for $n = 1, \ldots, N_\mathrm{times}$:
        \begin{enumerate}
          \item \textbf{Build NL self-collision term}: call
                \texttt{main\_nlterm}$(f^n)$ to compute the self-collision
                coefficient arrays \texttt{sc**} from the current distribution
                via the phi-distance kernel.
          \item \textbf{Update Coulomb logarithm} and rebuild the linear FP
                coefficients.
          \item \textbf{Assemble total operator}: $\Lcal = \Lcal_\mathrm{lin} + \Lcal_\mathrm{SC}$.
          \item \textbf{Rebuild LHS}: $M_\mathrm{lhs} = I - \theta\,\dt\,\Lcal$.
          \item \textbf{Compute RHS}:
                $b = (I+(1-\theta)\,\dt\,\Lcal)\,f^n + \dt\,S$.
          \item \textbf{Solve} (PARDISO phases 22+33 every step since $M_\mathrm{lhs}$ changes).
          \item Compute diagnostics.
        \end{enumerate}
  \item \textbf{Finalize} PARDISO.
\end{enumerate}

\textit{Note on start-up:} when starting from $f = 0$ with a beam source
(i.e., \texttt{istart = 0}), the NL self-collision term is skipped for the
first few steps until the beam density becomes significant enough for the
Rosenbluth potentials to be numerically meaningful.

\subsection{Multi-Phase Time-Stepping}

The namelist parameters \texttt{ntimes} and \texttt{timestep} are each
3-element arrays, enabling up to three sequential phases with different time
steps:
\begin{itemize}
  \item Phase~$k$ runs \texttt{ntimes(k)} steps of duration \texttt{timestep(k)} seconds.
  \item Phases with \texttt{ntimes(k) = 0} are skipped.
  \item Phases 2 and 3 append to the output files (restart in time).
\end{itemize}
Old single-phase input files (\texttt{ntimes = N, timestep = dt}) remain
backward-compatible: the scalar value fills element~1 and elements 2--3
default to zero.

\subsection{Initial Condition for Time-Dependent Runs}

The initial distribution $f^0$ is selected by \texttt{istart}:

\begin{center}
\begin{tabular}{cl}
\toprule
\texttt{istart} & Initial condition \\
\midrule
0 & $f^0 = 0$ everywhere. Requires \texttt{isource = -1}. \\
1 & Stix Maxwellian solution (default). \\
2 & Steady-state solution of the linear solver without self-collisions. \\
3 & Steady-state solution of the linear solver with Maxwellian SC background. \\
\bottomrule
\end{tabular}
\end{center}

\subsection{Steady-State Convergence Detection}

When \texttt{i\_ss\_check = -1}, the code monitors the relative \emph{rate of
change} of four indicators over a rolling window of \texttt{n\_ss\_window}
steps: the total kinetic energy $\Delta E/E$, the perpendicular kinetic energy
$\Delta E_\perp/E_\perp$ (which equilibrates last under RF drive and is the most
discriminating), the RF power $\Delta P_\mathrm{RF}/P_\mathrm{drive}$, and the
net power $\Delta P/P_\mathrm{drive}$.  Each relative change is divided by the
\emph{physical duration} of the window ($t_\mathrm{now}-t_{\mathrm{now}-W}$), so
the criterion measures a relative change \emph{per second}.  When all four rates
fall below \texttt{ss\_tol}, the run stops early with a ``CONVERGED'' message.

Because the window duration is taken from the simulation time rather than a step
count, \texttt{ss\_tol} is \emph{independent of the time step}: two runs with
different \texttt{timestep} that are equally close to steady state report the
same rate, so the same \texttt{ss\_tol} enforces the same physical stringency
(the duration is also measured correctly across phase boundaries where
\texttt{timestep} changes).  \textbf{Units:} \texttt{ss\_tol} is a per-second
relative rate; e.g.\ \texttt{ss\_tol = 3.d-4} stops the run when all four
indicators change by less than $3\times10^{-4}$ per second of simulated time.

\subsection{Diagnostics and Output}

At each time step the following quantities are written to disk:
\begin{itemize}
  \item Density $n(t) = \int f\,J\,\mathrm{d}^3v$.
  \item Total, perpendicular, and parallel kinetic energies.
  \item Effective temperature $T_\mathrm{eff} = (2T_\perp + T_\parallel)/3$,
        where $T_\perp = \tfrac{1}{2}m\langle \vp^2 \rangle$ and
        $T_\parallel = m\langle \vpa^2 \rangle$.
  \item Density-characteristic temperature $T_n$ (the cold-bulk temperature from
        the $\ln f$-vs-$v^2$ log-slope fit; it is the background used by
        \texttt{isc = 3} but is written as a diagnostic for \emph{all} runs).
  \item Near-axis and whole-grid minima of $f$ (a negative-$f$ monitor: small
        negative undershoots near the axis can corrupt low-velocity moments such
        as $T_n$).
  \item Perpendicular anisotropy factor $= 100\,T_\perp/(T_\perp + T_\parallel)$
        (\%), ranging from 0\% (all-parallel) to 50\% (Maxwellian) to 100\%
        (all-perpendicular).
  \item Power densities (collisional by species, RF, beam source, losses,
        self-collision total), and their perpendicular/parallel split for
        self-collisions.
  \item Momentum transfer rates (perpendicular and parallel) for each
        operator.
  \item Coulomb logarithm(s) vs.\ time.
\end{itemize}

At the end of a run, balance diagnostics (power, density, and optionally
momentum) are printed to standard output.

% ============================================================
\section{Namelist Reference}
\label{sec:namelist}
% ============================================================

The code reads its input from standard input (\texttt{stdin}) as a Fortran
\texttt{NAMELIST} group named \texttt{INPUT}.  A typical invocation is:
\begin{lstlisting}
FP2D_QLRF_NL.exe < jet_beam_case1.dat
\end{lstlisting}
The namelist file may omit any variable that has a default value; the code
will use the default.

\bigskip
The parameters are grouped thematically below.

% -------------------------------------------------------------------
\subsection{General and Output Control}

\begin{longtable}{>{\ttfamily}p{3.2cm} p{1.5cm} p{9cm}}
\toprule
\normalfont Parameter & Default & Description \\
\midrule
\endhead
\bottomrule
\endlastfoot

casename & \texttt{''} &
  Character string (max.\ 64 chars) appended to every output filename and
  plot title, e.g.\ \texttt{'JET-beam7-TD0-NLSC'}.  Leave empty for
  unnamed runs. \\[4pt]

iplot\_pow & $-1$ &
  Controls writing of power-vs-time output files.
  $-1$: write all power files.
  $0$: skip (useful when only energy and density diagnostics are needed). \\[4pt]

iplot\_mom & $0$ &
  Controls writing of momentum-transfer-rate files.
  $-1$: write all momentum files.
  $0$: skip (default; momentum files are large). \\[4pt]

idiag & $0$ &
  Selects the end-of-run balance diagnostics printed to stdout.
  $0$: power balance only.
  $-1$: all three balances (particle density, momentum, power). \\[4pt]

notxt & $0$ &
  Controls how many \texttt{.txt} output files are written (the \texttt{.dat}
  restart/kernel files are never affected).
  $0$: write all \texttt{.txt} outputs (normal).
  $1$: suppress all \texttt{.txt} outputs.
  $2$: write only the minimal test/diagnostic set, namely
  \texttt{Ekin\_perp}, \texttt{Ekin\_perp\_at\_vpar0}, \texttt{fout\_at\_vpar0},
  \texttt{Teff\_vs\_time}, and \texttt{anisotropy\_vs\_time}.
  Suppressed files are routed to the OS null device, so the run is otherwise
  unchanged (all stdout diagnostics still print). \\[4pt]

\end{longtable}

% -------------------------------------------------------------------
\subsection{Velocity-Space Grid}

\begin{longtable}{>{\ttfamily}p{3.2cm} p{1.5cm} p{9cm}}
\toprule
\normalfont Parameter & Default & Description \\
\midrule
\endhead
\bottomrule
\endlastfoot

nperp & -- &
  Number of grid points in $\vp$ (perpendicular direction). \\[4pt]

npar & -- &
  Number of grid points in $\vpa$ (parallel direction). \\[4pt]

vperp\_min & -- &
  Minimum perpendicular velocity (m/s).  Should be 0 or very small. \\[4pt]

vperp\_max & -- &
  Maximum perpendicular velocity (m/s).  Should be at least $3$--$5\times$
  the beam or thermal velocity. \\[4pt]

vpar\_min & -- &
  Minimum parallel velocity (m/s).  Must be negative. \\[4pt]

vpar\_max & -- &
  Maximum parallel velocity (m/s).  Must be positive and equal to
  $|$\texttt{vpar\_min}$|$ for symmetric grids. \\[4pt]

ising & $0$ &
  Grid spacing mode in $\vp$.
  $0$: uniform.
  $-1$: two-domain (fine for $\vp \le \texttt{vbound}$, coarser above).
  $+1$: quadratic spacing, higher resolution near $\vp = 0$. \\[4pt]

nsing & -- &
  Number of fine grid points in the inner domain when \texttt{ising = -1}.
  Must satisfy $1 \le \texttt{nsing} < \texttt{nperp}$. \\[4pt]

vbound & -- &
  Boundary between fine and coarse regions in $\vp$ when \texttt{ising = -1}
  (m/s). \\[4pt]

\end{longtable}

% -------------------------------------------------------------------
\subsection{Plasma Parameters}

\begin{longtable}{>{\ttfamily}p{3.2cm} p{1.5cm} p{9cm}}
\toprule
\normalfont Parameter & Default & Description \\
\midrule
\endhead
\bottomrule
\endlastfoot

nbulk & -- &
  Number of background species (including electrons).
  Species 1 is always electrons.  Up to 10 species are supported. \\[4pt]

t(1:nbulk) & -- &
  Temperatures of background species (eV).
  \texttt{t(1)}: electron temperature $T_e$.
  \texttt{t(2)}: temperature of first ion species, etc. \\[4pt]

aa & -- &
  Mass number of the minority (test) ion species. \\[4pt]

za & -- &
  Charge number (atomic number) of the minority ion species. \\[4pt]

ab(1:9) & -- &
  Mass numbers of background ion species 2, 3, \ldots\ (species 1 is
  always electrons with $m_e/m_p$). \\[4pt]

zb(1:9) & -- &
  Charge numbers of background ion species 2, 3, \ldots \\[4pt]

ne & -- &
  Electron density (m$^{-3}$). \\[4pt]

xpart & -- &
  Concentration of minority ion species relative to $n_e$,
  i.e.\ $n_a = \texttt{xpart} \times n_e$. \\[4pt]

xb(1:9) & -- &
  Concentrations of background ion species 2, 3, \ldots\ relative to $n_e$.
  Quasi-neutrality is assumed: $\sum_b Z_b\,n_b = n_e$. \\[4pt]

\end{longtable}

% -------------------------------------------------------------------
\subsection{Beam Source}

\begin{longtable}{>{\ttfamily}p{3.2cm} p{1.5cm} p{9cm}}
\toprule
\normalfont Parameter & Default & Description \\
\midrule
\endhead
\bottomrule
\endlastfoot

isource & -- &
  Beam source flag.
  $-1$: include NBI beam source and particle losses.
  $0$: no source; distribution is sourceless (total particles conserved). \\[4pt]

beam\_ekin & -- &
  Beam injection energy (keV). \\[4pt]

beam\_angle\_deg & -- &
  Injection angle (degrees) measured from the magnetic field direction.
  $0^\circ$: purely parallel injection.
  $90^\circ$: purely perpendicular injection. \\[4pt]

beam\_dvperp & -- &
  Width $\sigma_\perp$ of the Gaussian approximation to the delta function
  in $\vp$ (m/s).  A value of $\sim 5$--$10\%$ of $\vp^\mathrm{beam}$ is
  typical; must be resolved by the grid. \\[4pt]

beam\_dvpar & -- &
  Width $\sigma_\parallel$ of the Gaussian in $\vpa$ (m/s). \\[4pt]

taus & -- &
  Slowing-down time $\tau_s$ (s).  Sets the source amplitude
  $S_0 = n_a / \tau_s$ and the loss rate. \\[4pt]

\end{longtable}

% -------------------------------------------------------------------
\subsection{RF Heating}

\begin{longtable}{>{\ttfamily}p{3.2cm} p{1.5cm} p{9cm}}
\toprule
\normalfont Parameter & Default & Description \\
\midrule
\endhead
\bottomrule
\endlastfoot

irf & -- &
  RF flag.
  $-1$: include quasi-linear RF operator.
  $0$ or else: no RF. \\[4pt]

frek & -- &
  Wave frequency $f$ (Hz). \\[4pt]

b0 & -- &
  Static magnetic field $B_0$ (T). \\[4pt]

nharm & -- &
  Resonance harmonic number $n$ (e.g.\ 1 for fundamental, 2 for second
  harmonic). \\[4pt]

kpar & -- &
  Parallel wave vector $k_\parallel$ (rad/m). \\[4pt]

kperp & -- &
  Perpendicular wave vector $k_\perp$ (complex, rad/m). \\[4pt]

eplus & -- &
  Left-hand polarisation amplitude $E_+$ (complex, V/m). \\[4pt]

emin & -- &
  Right-hand polarisation amplitude $E_-$ (complex, V/m). \\[4pt]

delta\_RF & -- &
  Width of the Gaussian broadening of the resonance in $\vpa$ (m/s).
  Must be wide enough to resolve the resonance on the grid. \\[4pt]

\end{longtable}

% -------------------------------------------------------------------
\subsection{Time Integration and Solver Control}

\begin{longtable}{>{\ttfamily}p{3.2cm} p{1.5cm} p{9cm}}
\toprule
\normalfont Parameter & Default & Description \\
\midrule
\endhead
\bottomrule
\endlastfoot

ntimes(1:3) & \texttt{0,0,0} &
  Number of time steps per phase.
  \texttt{ntimes(1) = 0}: steady-state run.
  \texttt{ntimes(1) > 0}: time-dependent run.
  Up to three phases with different step sizes are supported; phases with
  \texttt{ntimes(k) = 0} are skipped.
  Old-format scalar values (\texttt{ntimes = N}) are backward-compatible and
  fill \texttt{ntimes(1)}. \\[4pt]

timestep(1:3) & \texttt{0,0,0} &
  Time step size per phase (s).
  \texttt{timestep(k)} is used for phase~$k$.
  Backward-compatible with scalar input. \\[4pt]

icn & -- &
  Time-differencing scheme.
  $-1$: Crank--Nicolson ($\theta = 0.5$).
  $+1$: intermediate scheme ($\theta = 0.75$).
  else: fully implicit ($\theta = 1.0$). \\[4pt]

iold & -- &
  Initial condition source.
  $0$: fresh start (use \texttt{istart} to select the initial condition).
  $-1$: restart from a previous run by reading \texttt{xout.dat}. \\[4pt]

istart & $1$ &
  Initial distribution for fresh time-dependent runs (\texttt{iold = 0}).
  $0$: $f^0 = 0$ (beam-driven only; requires \texttt{isource = -1}).
  $1$: Stix Maxwellian.
  $2$: Steady-state solution without self-collisions.
  $3$: Steady-state solution with Maxwellian self-collisions. \\[4pt]

isc & -- &
  Self-collision treatment.
  $0$: no self-collisions (linear FP equation).
  $1$: Maxwellian background at fixed Stix temperature.
  $2$: Maxwellian background at evolving effective temperature $T_\mathrm{eff}(t)$
       (requires TD run).
  $3$: Maxwellian background at the density-characteristic temperature $T_n(t)$
       (requires TD run; see \texttt{core\_frac}).
  $-1$: fully nonlinear self-collision operator (requires TD run). \\[4pt]

core\_frac & $3.8\!\times\!10^{-3}$ &
  Core fraction defining the $T_n$ log-slope fit for \texttt{isc = 3}: only grid
  cells with $f > \texttt{core\_frac}\cdot\max f$ enter the $\ln f$-vs-$v^2$ fit.
  Smaller values weight the cold bulk more heavily.  Ignored unless
  \texttt{isc = 3}; the appropriate value is scenario-dependent. \\[4pt]

i\_upwind & $0$ &
  Advection stabilization of the perpendicular drag term $B\,\partial f/\partial\vp$
  (Section~\ref{sec:upwind}).
  $0$: central Fornberg weights (default).
  $-1$: P\'eclet-hybrid upwinding --- in cells where the cell-P\'eclet number
  $|B|\,\Delta\vp/D > 2$ the drag term uses a one-sided (first-order) upwind
  difference, restoring dissipativity of the advection-dominated outer boundary
  layer and removing the Crank--Nicolson high-$\vp$ instability.
  $1$: P\'eclet-hybrid \emph{power-law} scheme (Patankar 1980) --- in the same
  advection-dominated cells the combined drag$+$diffusion operator is
  discretised with a 3-point control-volume power-law stencil, which removes the
  instability with roughly half the tail numerical diffusion of the first-order
  upwind ($-1$).  Both leave the high-order 7-point Fornberg stencil in the
  smooth bulk. \\[4pt]

\end{longtable}

% -------------------------------------------------------------------
\subsection{Steady-State Convergence Detection}

\begin{longtable}{>{\ttfamily}p{3.2cm} p{1.5cm} p{9cm}}
\toprule
\normalfont Parameter & Default & Description \\
\midrule
\endhead
\bottomrule
\endlastfoot

i\_ss\_check & $0$ &
  Enable automatic convergence detection.
  $0$: disabled (run for the full \texttt{ntimes} steps).
  $-1$: stop early when all convergence criteria are met. \\[4pt]

n\_ss\_window & $50$ &
  Rolling window width (number of time steps) over which the
  convergence indicators $\Delta E/E$, $\Delta n/n$, $\Delta P/P$ are
  evaluated. \\[4pt]

ss\_tol & $10^{-3}$ &
  Relative convergence tolerance.  All three indicators must fall below
  this value for the run to stop. \\[4pt]

\end{longtable}


% ============================================================
\section{Output Files}
% ============================================================

All output files are written to the working directory.  When \texttt{casename}
is non-empty, the string \texttt{-casename} is inserted before the file
extension.

\subsection{Final Distribution}

\begin{description}
  \item[\texttt{fout.txt}] Final VDF as a list of $(\vp, \vpa, f)$ triples.
  \item[\texttt{xout.dat}] Binary restart file: one header line (final time)
        followed by the packed 1D solution vector.  Used with \texttt{iold = -1}.
  \item[\texttt{fstix.dat}] Stix Maxwellian reference distribution.
\end{description}

\subsection{1D Cuts of the VDF}

\texttt{fout\_at\_vpar0.txt}, \texttt{fout\_at\_vperp0.txt},
\texttt{fout\_at\_vperpmax.txt}, \texttt{fout\_at\_vparmax.txt}:
values of $f$ along the four boundary slices of the domain.

\subsection{Time Traces}

Written at every time step during the simulation:

\begin{description}
  \item[\texttt{density\_vs\_time.txt}] Columns: time (s), $n(t)$ (m$^{-3}$).
  \item[\texttt{energy\_vs\_time.txt}] Columns: time, $E_\mathrm{tot}$ (keV), $E_\perp$ (keV).
  \item[\texttt{anisotropy\_vs\_time.txt}] Columns: time, perpendicular anisotropy (\%).
  \item[\texttt{Teff\_vs\_time.txt}] Columns: time, $T_\mathrm{eff}$ (keV).
  \item[\texttt{Tn\_vs\_time.txt}] Columns: time, $T_n$ (keV) -- density-characteristic
        (cold-bulk) temperature; written for all runs.
  \item[\texttt{power\_coll\_tot\_vs\_time.txt}] Total collisional power density (MW/m$^3$).
  \item[\texttt{power\_coll\_e\_vs\_time.txt}] Electron collisional power density.
  \item[\texttt{power\_coll\_ion$N$\_vs\_time.txt}] Ion $N$ collisional power density.
  \item[\texttt{power\_coll\_self\_vs\_time.txt}] Self-collision power density
        (4 columns: time, total, $P_\mathrm{SC}^\perp$, $P_\mathrm{SC}^\parallel$).
  \item[\texttt{power\_NBI\_vs\_time.txt}] Beam injection and particle loss power.
  \item[\texttt{power\_RF\_vs\_time.txt}] RF power density.
  \item[\texttt{coulomb\_log\_vs\_time.txt}] Coulomb logarithm for background ion species.
  \item[\texttt{coulomb\_log\_self\_vs\_time.txt}] Self-collision Coulomb logarithm.
  \item[\texttt{momentum\_coll\_*\_vs\_time.txt}] Momentum transfer rates
        (2 columns per file: $\perp$ and $\parallel$ components), written when
        \texttt{iplot\_mom = -1}.
\end{description}

\subsection{End-of-Run Balance Diagnostics (stdout)}

A power balance (and, when \texttt{idiag = -1}, particle density and
momentum balances) is printed to standard output at the end of the run.
Values are given in MW/m$^3$ for power, N/m$^3$ for momentum, and s$^{-1}$
for particle rates.


% ============================================================
\section{Example Namelist}
% ============================================================

The following example illustrates a time-dependent D-beam simulation in a D
plasma with nonlinear self-collisions, using a two-phase time-stepping
strategy: 1000 coarse steps followed by 1000 finer steps. It lists \emph{every}
\texttt{\&INPUT} parameter; the RF block is included for completeness but is only
read when \texttt{irf=-1}, and the diagnostic/advanced flags are shown at their
default (off) values.

\begin{lstlisting}
&INPUT
  casename = 'JET-beam7-TD0-NLSC',

  ! Grid
  nperp = 151, npar = 151,
  vperp_min = 0.d0,  vperp_max = 5.d6,
  vpar_min  = -4.d6, vpar_max  = 4.d6,
  ising = 1, nsing = 25, vbound = 0.07d6,

  ! Plasma: electrons + D ions
  nbulk = 2,
  t     = 15.d3, 10.d3, 8*0.d0,
  aa = 2.d0, za = 1.d0,
  ab = 2.d0, 8*0.d0,
  zb = 1.d0, 8*0.d0,
  ne     = 5.0d19,
  xpart  = 0.05d0,
  xb     = 0.50d0, 8*0.d0,

  ! D-NBI beam at 120 keV, 90 degrees
  isource = -1,
  beam_ekin      = 120.d0,
  beam_angle_deg = 90.d0,
  beam_dvperp    = 0.12d6,
  beam_dvpar     = 0.20d6,
  taus = 300.d-3,

  ! RF quasi-linear operator (b0..delta_RF read only when irf = -1)
  irf   = 0,
  b0    = 3.45d0,
  frek  = 53.d6,
  nharm = 1,
  kpar  = 2.d0,
  kperp = (30.d0, 1.d0),
  eplus = (1.d3, 3.25d3),
  emin  = (3.d3, 3.75d3),
  delta_RF = 0.5d6,

  ! Time integration: two phases
  icn      = 1,
  ntimes   = 1000, 1000, 0,
  timestep = 1.d-3, 10.d-3, 0.d0,
  iold   = 0,
  istart = 0,
  isc    = -1,
  core_frac = 3.8d-3,

  ! Steady-state convergence detection
  i_ss_check  = -1,
  n_ss_window = 50,
  ss_tol      = 1.d-4,

  ! Output
  iplot_pow = -1,
  iplot_mom = -1,
  idiag     = -1,
  notxt     = 0,

  ! Advanced numerics (default 0 = off)
  i_upwind     = 0,
/
\end{lstlisting}


% ============================================================
\section{Compilation and Execution}
% ============================================================

\subsection{Build}

Open the Visual Studio solution
\texttt{FP2D\_QLRF\_NL/FP2D\_QLRF\_NL.sln} and build the
\texttt{Release|x64} configuration.  The compiler is Intel \texttt{ifx}
with Intel MKL (sequential).  The executable is placed at
\texttt{FP2D\_QLRF\_NL/x64/Release/FP2D\_QLRF\_NL.exe}.

\subsection{Running}

The program reads from \texttt{stdin}:
\begin{lstlisting}
FP2D_QLRF_NL.exe < inputs/jet_beam_case1.dat
\end{lstlisting}
Output files are written to the working directory.

\subsection{Plotting}

Use the companion Python script \texttt{fp2d\_plot.py}:
\begin{lstlisting}
python fp2d_plot.py run FP2D_QLRF_NL.exe inputs/jet_beam_case1.dat --show
python fp2d_plot.py plot x64/Release --casename JET-beam7-TD0-NLSC --save plots/
python fp2d_plot.py compare x64/Release --cases JET-beam7-TD0-Lin JET-beam7-TD0-NLSC
\end{lstlisting}
See \texttt{fp2d\_plot-readme.txt} for the full documentation of the
plotting tool.

\end{document}
